% theme: trigonometric_equations
\MajorQuestionLesson{三角方程式}
\SetRowSpace{4mm}

\TypeQuestionSingle{次の方程式を、$0\leqq\theta<2\pi$ の範囲で解きなさい。}{basic_value}
\QQRowThree{$\sin\theta=\dfrac{1}{2}$}{$\dfrac{\pi}{6},\ \dfrac{5\pi}{6}$}{$\cos\theta=-\dfrac{1}{2}$}{$\dfrac{2\pi}{3},\ \dfrac{4\pi}{3}$}{$\tan\theta=1$}{$\dfrac{\pi}{4},\ \dfrac{5\pi}{4}$}
\QQRowThree{$\cos\theta=\dfrac{\sqrt{2}}{2}$}{$\dfrac{\pi}{4},\ \dfrac{7\pi}{4}$}{$\tan\theta=-\sqrt{3}$}{$\dfrac{2\pi}{3},\ \dfrac{5\pi}{3}$}{$\sin\theta=-\dfrac{\sqrt{2}}{2}$}{$\dfrac{5\pi}{4},\ \dfrac{7\pi}{4}$}
\QQRowThree{$\tan\theta=0$}{$0,\ \pi$}{$\sin\theta=1$}{$\dfrac{\pi}{2}$}{$\cos\theta=-\dfrac{\sqrt{3}}{2}$}{$\dfrac{5\pi}{6},\ \dfrac{7\pi}{6}$}
\QQRowThree{$\sin\theta=-\dfrac{\sqrt{3}}{2}$}{$\dfrac{4\pi}{3},\ \dfrac{5\pi}{3}$}{$\cos\theta=0$}{$\dfrac{\pi}{2},\ \dfrac{3\pi}{2}$}{$\tan\theta=\dfrac{\sqrt{3}}{3}$}{$\dfrac{\pi}{6},\ \dfrac{7\pi}{6}$}
\QQRowThree{$\cos\theta=-1$}{$\pi$}{$\tan\theta=-1$}{$\dfrac{3\pi}{4},\ \dfrac{7\pi}{4}$}{$\sin\theta=0$}{$0,\ \pi$}
\QQRowThree{$\tan\theta=\sqrt{3}$}{$\dfrac{\pi}{3},\ \dfrac{4\pi}{3}$}{$\sin\theta=\dfrac{\sqrt{3}}{2}$}{$\dfrac{\pi}{3},\ \dfrac{2\pi}{3}$}{$\cos\theta=\dfrac{\sqrt{3}}{2}$}{$\dfrac{\pi}{6},\ \dfrac{11\pi}{6}$}

\TypeQuestionSingle{次の方程式を、$0\leqq\theta<2\pi$ の範囲で解きなさい。}{linear_form}
\QQRow{$2\sin\theta-\sqrt{3}=0$}{$\dfrac{\pi}{3},\ \dfrac{2\pi}{3}$}{$\sqrt{2}\cos\theta+1=0$}{$\dfrac{3\pi}{4},\ \dfrac{5\pi}{4}$}
\QQRow{$3\tan\theta+\sqrt{3}=0$}{$\dfrac{5\pi}{6},\ \dfrac{11\pi}{6}$}{$4\cos\theta-2=0$}{$\dfrac{\pi}{3},\ \dfrac{5\pi}{3}$}
\QQRow{$2\sin\theta+1=0$}{$\dfrac{7\pi}{6},\ \dfrac{11\pi}{6}$}{$\sqrt{2}\sin\theta-1=0$}{$\dfrac{\pi}{4},\ \dfrac{3\pi}{4}$}
\QQRow{$\sqrt{3}\tan\theta-1=0$}{$\dfrac{\pi}{6},\ \dfrac{7\pi}{6}$}{$2\tan\theta-2\sqrt{3}=0$}{$\dfrac{\pi}{3},\ \dfrac{4\pi}{3}$}

\AnswerColumnBreak
\TypeQuestionSingle{次の方程式を、$0\leqq\theta<2\pi$ の範囲で解きなさい。}{multiple_angle}
\QQRow{$\sin2\theta=\dfrac{\sqrt{3}}{2}$}{\AnswerTwoLines{$\dfrac{\pi}{6},\ \dfrac{\pi}{3}$}{$\dfrac{7\pi}{6},\ \dfrac{4\pi}{3}$}}{$\cos2\theta=-\dfrac{\sqrt{2}}{2}$}{\AnswerTwoLines{$\dfrac{3\pi}{8},\ \dfrac{5\pi}{8}$}{$\dfrac{11\pi}{8},\ \dfrac{13\pi}{8}$}}
\QQRow{$\tan2\theta=1$}{\AnswerTwoLines{$\dfrac{\pi}{8},\ \dfrac{5\pi}{8}$}{$\dfrac{9\pi}{8},\ \dfrac{13\pi}{8}$}}{$\sin3\theta=-1$}{$\dfrac{\pi}{2},\ \dfrac{7\pi}{6},\ \dfrac{11\pi}{6}$}
\QQRow{$\cos\left(\theta-\dfrac{\pi}{3}\right)=\dfrac{1}{2}$}{$0,\ \dfrac{2\pi}{3}$}{$\tan\left(\theta+\dfrac{\pi}{6}\right)=-\sqrt{3}$}{$\dfrac{\pi}{2},\ \dfrac{3\pi}{2}$}

\LessonPageBreak

\TypeQuestionSingle{次の方程式を、$0\leqq\theta<2\pi$ の範囲で解きなさい。}{quadratic_one_trig}
\QQRow{$2\sin^2\theta-3\sin\theta+1=0$}{$\dfrac{\pi}{6},\ \dfrac{\pi}{2},\ \dfrac{5\pi}{6}$}{$2\cos^2\theta+\cos\theta-1=0$}{$\dfrac{\pi}{3},\ \pi,\ \dfrac{5\pi}{3}$}
\QQRow{$\tan^2\theta-\tan\theta=0$}{\AnswerTwoLines{$0,\ \dfrac{\pi}{4}$}{$\pi,\ \dfrac{5\pi}{4}$}}{$4\sin^2\theta-1=0$}{\AnswerTwoLines{$\dfrac{\pi}{6},\ \dfrac{5\pi}{6}$}{$\dfrac{7\pi}{6},\ \dfrac{11\pi}{6}$}}
\QQRow{$2\cos^2\theta-3\cos\theta=0$}{$\dfrac{\pi}{2},\ \dfrac{3\pi}{2}$}{$\tan^2\theta-3=0$}{\AnswerTwoLines{$\dfrac{\pi}{3},\ \dfrac{2\pi}{3}$}{$\dfrac{4\pi}{3},\ \dfrac{5\pi}{3}$}}
\QQRow{$2\sin^2\theta+\sin\theta-1=0$}{$\dfrac{\pi}{6},\ \dfrac{5\pi}{6},\ \dfrac{3\pi}{2}$}{$4\cos^2\theta-3=0$}{\AnswerTwoLines{$\dfrac{\pi}{6},\ \dfrac{5\pi}{6}$}{$\dfrac{7\pi}{6},\ \dfrac{11\pi}{6}$}}

\TypeQuestionSingle{三角関数の相互関係を用いて、次の方程式を $0\leqq\theta<2\pi$ の範囲で解きなさい。}{identity_equation}
\QQRow{$2\sin^2\theta+3\cos\theta-3=0$}{$0,\ \dfrac{\pi}{3},\ \dfrac{5\pi}{3}$}{$2\cos^2\theta+\sin\theta-1=0$}{$\dfrac{\pi}{2},\ \dfrac{7\pi}{6},\ \dfrac{11\pi}{6}$}
\QQRow{$\sin^2\theta=3\cos^2\theta$}{\AnswerTwoLines{$\dfrac{\pi}{3},\ \dfrac{2\pi}{3}$}{$\dfrac{4\pi}{3},\ \dfrac{5\pi}{3}$}}{$2\sin^2\theta-\cos^2\theta=\dfrac{1}{2}$}{\AnswerTwoLines{$\dfrac{\pi}{4},\ \dfrac{3\pi}{4}$}{$\dfrac{5\pi}{4},\ \dfrac{7\pi}{4}$}}
\QQRow{$3\sin^2\theta+\cos^2\theta=2$}{\AnswerTwoLines{$\dfrac{\pi}{4},\ \dfrac{3\pi}{4}$}{$\dfrac{5\pi}{4},\ \dfrac{7\pi}{4}$}}{$\sin^2\theta+\cos^2\theta+\tan^2\theta=4$}{\AnswerTwoLines{$\dfrac{\pi}{3},\ \dfrac{2\pi}{3}$}{$\dfrac{4\pi}{3},\ \dfrac{5\pi}{3}$}}
\QQRow{$\dfrac{1-\sin^2\theta}{1+\tan^2\theta}=\dfrac{1}{4}$}{\AnswerTwoLines{$\dfrac{\pi}{4},\ \dfrac{3\pi}{4}$}{$\dfrac{5\pi}{4},\ \dfrac{7\pi}{4}$}}{$\dfrac{\sin^2\theta}{\cos^2\theta}=\dfrac{1}{3}$}{\AnswerTwoLines{$\dfrac{\pi}{6},\ \dfrac{5\pi}{6}$}{$\dfrac{7\pi}{6},\ \dfrac{11\pi}{6}$}}

\AnswerColumnBreak
\TypeQuestionSingle{次の方程式を、$0\leqq\theta<2\pi$ の範囲で解きなさい。}{mixed_and_product}
\QQRow{$\sin\theta\left(2\cos\theta-1\right)=0$}{$0,\ \dfrac{\pi}{3},\ \pi,\ \dfrac{5\pi}{3}$}{$\cos\theta\left(2\sin\theta+1\right)=0$}{\AnswerTwoLines{$\dfrac{\pi}{2},\ \dfrac{7\pi}{6}$}{$\dfrac{3\pi}{2},\ \dfrac{11\pi}{6}$}}
\QQRow{$\sin\theta=\cos\theta$}{$\dfrac{\pi}{4},\ \dfrac{5\pi}{4}$}{$\sqrt{3}\sin\theta+\cos\theta=0$}{$\dfrac{5\pi}{6},\ \dfrac{11\pi}{6}$}
\QQRow{$\left(2\sin\theta-1\right)\left(2\cos\theta+\sqrt{2}\right)=0$}{\AnswerTwoLines{$\dfrac{\pi}{6},\ \dfrac{3\pi}{4}$}{$\dfrac{5\pi}{6},\ \dfrac{5\pi}{4}$}}{$\sin^2\theta-\cos^2\theta=0$}{\AnswerTwoLines{$\dfrac{\pi}{4},\ \dfrac{3\pi}{4}$}{$\dfrac{5\pi}{4},\ \dfrac{7\pi}{4}$}}
\QQRow{$\sin\theta+\cos\theta=0$}{$\dfrac{3\pi}{4},\ \dfrac{7\pi}{4}$}{$\sin\theta-\sqrt{3}\cos\theta=0$}{$\dfrac{\pi}{3},\ \dfrac{4\pi}{3}$}

\AnswerColumnBreak
\TypeQuestionDouble{$n$ を整数として、次の方程式の一般解を求めなさい。}{general_solution}
\QQRowThree{$\sin\theta=-\dfrac{1}{2}$}{\AnswerTwoLines{$\dfrac{7\pi}{6}+2n\pi,$}{$\dfrac{11\pi}{6}+2n\pi$}}{$\cos\theta=\dfrac{\sqrt{2}}{2}$}{$2n\pi\pm\dfrac{\pi}{4}$}{$\tan\theta=-\sqrt{3}$}{$-\dfrac{\pi}{3}+n\pi$}
\QQRowThree{$\sin2\theta=\dfrac{\sqrt{3}}{2}$}{\AnswerTwoLines{$\dfrac{\pi}{6}+n\pi,$}{$\dfrac{\pi}{3}+n\pi$}}{$\cos\left(\theta+\dfrac{\pi}{4}\right)=0$}{$\dfrac{\pi}{4}+n\pi$}{$2\sin\theta-\sqrt{3}=0$}{\AnswerTwoLines{$\dfrac{\pi}{3}+2n\pi,$}{$\dfrac{2\pi}{3}+2n\pi$}}
\QQRowThree{$\sqrt{2}\cos\theta+1=0$}{\AnswerTwoLines{$\dfrac{3\pi}{4}+2n\pi,$}{$\dfrac{5\pi}{4}+2n\pi$}}{$3\tan\theta+\sqrt{3}=0$}{$-\dfrac{\pi}{6}+n\pi$}{$\tan\left(2\theta-\dfrac{\pi}{4}\right)=1$}{$\dfrac{\pi}{4}+\dfrac{n\pi}{2}$}
